\section{Exact coordinate permutations at two scales}\label{pairs:section}

The tensor condition will now produce two configurations with the same
coordinate values at each label, but with arbitrarily different scales.
The equality of coordinate values is an equality of finite multisets,
including multiplicities.

\begin{lemma}[Two-scale configurations]\label{pairs:two-scale}
Let $A=\{a_1,\ldots,a_s\}\subset\R^d$ have affine span $\R^d$, with
$s\geq2$. Let $W=\R^{1+d}$ be regarded as a vector space over $\Q$.
Suppose that a symmetric tensor $T\in W\otimes_{\Q}W$ satisfies
\[
 (e_i\otimes e_i)(T)=0\quad(1\leq i\leq s),
 \qquad G(T)=I_d,
\]
where $e_i(v,u)=v+a_i\cdot u$. For every $\eta>0$, there are an integer
$\ell\geq1$ and vectors $f_i,g_i\in\R^\ell$ such that
\begin{align}
 \|f_i-f_j\|^2&=\|a_i-a_j\|^2,
 &\|g_i-g_j\|^2&=q\|a_i-a_j\|^2,
 &q&=\frac{\eta}{1+\eta}, \label{pairs:distances}
\end{align}
for every $i,j$, and $g_i$ is a permutation of the coordinates of $f_i$
for each $i$.
\end{lemma}

The finite function $\nu$ from Lemma~\ref{stencil:finite} supplies signed rational
weights whose sums vanish on every evaluation class. Positive and negative
weights can therefore become two lists of coordinates with identical
multisets of values. Their second moments control the difference of their
gradient Gram matrices. The remaining task is to make the negative list's
Gram matrix arbitrarily small. We first obtain the required inequality
by smoothing, and then obtain finite rational weights while preserving
all cancellation identities exactly.

For a real symmetric $k\times k$ matrix $C$, write
\[
 D_C=\frac12\sum_{j,l=1}^k C_{jl}\partial_j\partial_l,
 \qquad r_+ = \max\{r,0\},\quad r_- = \max\{-r,0\}\quad(r\in\R).
\]

\begin{lemma}[A compactly supported smoothing function]\label{pairs:bump}
Let $C$ be a real symmetric $k\times k$ matrix and let
$U\in\operatorname{Mat}_{d\times k}(\R)$, where $d\geq1$ and
$UCU^{\mathsf T}=I_d$. For every $\eta>0$, there is a nonnegative
$b\in C_c^\infty(\R^k)$ such that
\begin{equation}\label{pairs:bump-bound}
 \int_{\R^k}b(x)\,dx>0,
 \qquad
 \int_{\R^k}(D_Cb)_-(x)\|Ux\|^2\,dx
 <\eta\int_{\R^k}b(x)\,dx.
\end{equation}
\end{lemma}

\begin{proof}
We will average Gaussian densities with covariance matrices running from
$S$ to $S+C$. We first construct a positive definite $S$ for which
$S+C$ is positive definite and $\tr(USU^{\mathsf T})$ is small.

The identity $UCU^{\mathsf T}=I_d$ implies that $U$ has rank $d$.
Set $H=\operatorname{im}(U^{\mathsf T})$ and $V=\ker U$, so that
$\R^k=H\oplus V$ orthogonally. The restriction of the quadratic form
$C$ to $H$ is positive definite: if $0\ne h=U^{\mathsf T}y$, then
\[
 h^{\mathsf T}Ch=y^{\mathsf T}UCU^{\mathsf T}y=\|y\|^2>0.
\]
Choose $\epsilon>0$ with
$\epsilon\tr(UU^{\mathsf T})<\eta/8$. We claim that, for a sufficiently
large $L>0$, the matrix
\[
 S=\epsilon P_H+L P_V
\]
is positive definite and so is $S+C$. Here $P_H,P_V$ are the orthogonal
projections. To check the latter assertion, write
\[
 S+C=
 \begin{pmatrix}
  A&B\\ B^{\mathsf T}&C_{VV}+L I_V
 \end{pmatrix},
 \qquad A=C_{HH}+\epsilon I_H>0.
\]
Choose $L$ such that
$C_{VV}+L I_V-B^{\mathsf T}A^{-1}B>0$ and use the Schur complement.
If $V=\{0\}$, the second block is absent. Since $UP_V=0$, this construction
also gives
\begin{equation}\label{pairs:small-covariance}
 \tr(USU^{\mathsf T})=\epsilon\tr(UU^{\mathsf T})<\eta/8.
\end{equation}
Every matrix $S+tC$, $0\leq t\leq1$, is positive definite, being a convex
combination of $S$ and $S+C$.

For a positive definite $k\times k$ matrix $Q$, let
\[
 \phi_Q(x)=(2\pi)^{-k/2}(\det Q)^{-1/2}
                 \exp\!\left(-\frac12x^{\mathsf T}Q^{-1}x\right)
\]
be the centered Gaussian density with covariance $Q$. Differentiating
the displayed formula, with $Q_t=S+tC$, gives
\[
 \frac{d}{dt}\phi_{Q_t}(x)
 =\frac12\left(
    x^{\mathsf T}Q_t^{-1}CQ_t^{-1}x-\tr(Q_t^{-1}C)
   \right)\phi_{Q_t}(x)
 =D_C\phi_{Q_t}(x).
\]
This identity does not require $C$ to be positive semidefinite. The
covariances $Q_t$ have uniformly bounded eigenvalues, bounded away from
zero, so the densities and each of their spatial derivatives have
uniform Gaussian decay. Consequently
\[
 b_*(x)=\int_0^1\phi_{S+tC}(x)\,dt
\]
is a nonnegative Schwartz function of integral $1$, and differentiation
under the integral gives
\[
 D_Cb_* = \phi_{S+C}-\phi_S.
\]
The pointwise inequality $(\phi_{S+C}-\phi_S)_-\leq\phi_S$ and
\eqref{pairs:small-covariance} imply
\begin{equation}\label{pairs:gaussian-bound}
 \int_{\R^k}(D_Cb_*)_-(x)\|Ux\|^2\,dx
 \leq\int_{\R^k}\phi_S(x)\|Ux\|^2\,dx
 =\tr(USU^{\mathsf T})<\eta/8.
\end{equation}

To make the support compact, choose a smooth function $\chi$ with
$0\leq\chi\leq1$, equal to $1$ on the unit ball and zero outside the
ball of radius $2$. Put $\chi_R(x)=\chi(x/R)$ and $b_R=\chi_Rb_*$.
Then $\int b_R\to1$. The product rule reads
\[
 D_Cb_R=\chi_R D_Cb_*
  +\sum_{j,l}C_{jl}(\partial_j\chi_R)(\partial_lb_*)
  +b_*D_C\chi_R.
\]
The derivatives of $\chi_R$ have bounds of orders $R^{-1}$ and $R^{-2}$
and are supported where $R\leq\|x\|\leq2R$. The Schwartz decay therefore
implies
\[
 \int_{\R^k}|D_Cb_R-D_Cb_*|(x)\|Ux\|^2\,dx\longrightarrow0.
\]
Since $r\mapsto r_-$ is $1$-Lipschitz, the corresponding weighted
negative-part integrals converge as well. For sufficiently large $R$,
we have $\int b_R>1/2$ and the last displayed integral is less than
$\eta/8$. Equation~\eqref{pairs:gaussian-bound} now gives
\[
 \int(D_Cb_R)_-\|Ux\|^2<\eta/4<\eta\int b_R.
\]
Taking $b=b_R$ proves the lemma.
\end{proof}

\begin{lemma}[Finite rational weights]\label{pairs:rational}
Let $\Gamma=\Z^k$, and let $\Lambda_1,\ldots,\Lambda_s$ be subgroups
of $\Gamma$. Suppose that a finitely supported function
$\nu:\Gamma\to\Q$ has zero pushforward on each $\Gamma/\Lambda_i$,
and satisfies
\[
 \sum_\gamma\nu(\gamma)=0,
 \qquad \sum_\gamma\nu(\gamma)\gamma=0,
 \qquad \sum_\gamma\nu(\gamma)\gamma\gamma^{\mathsf T}=C.
\]
Let $U$ be a real $d\times k$ matrix with $UCU^{\mathsf T}=I_d$,
where $d\geq1$. For every $\eta>0$, there are an integer $n\geq1$,
a finitely supported $\mu:\Gamma\to\Q$, and $\zeta\in\Q_{>0}$ such
that $\mu$ has zero pushforward on each $\Gamma/\Lambda_i$ and
\begin{align}
 \sum_\lambda\mu(\lambda)
       \frac{\lambda}{n}\frac{\lambda^{\mathsf T}}{n}
   &=\zeta C, \label{pairs:exact-moment}\\
 \sum_\lambda\mu(\lambda)_-
       \left\|U\frac{\lambda}{n}\right\|^2
   &<\eta\zeta. \label{pairs:small-negative}
\end{align}
\end{lemma}

\begin{proof}
Take $b$ from Lemma~\ref{pairs:bump}. For each positive integer $n$, set
\[
 d_n(\lambda)=n^{2-k}b(\lambda/n),
 \qquad
 \mu_n(\lambda)=(\nu*d_n)(\lambda)
      =\sum_\gamma\nu(\gamma)d_n(\lambda-\gamma).
\]
These functions have finite support. Quotient pushforward commutes with
convolution, so every required pushforward of $\mu_n$ is zero.
More generally, for any finitely supported real function $d$ on $\Gamma$,
the moment identities for $\nu$ give
\begin{equation}\label{pairs:convolution-moments}
 \sum_\lambda(\nu*d)(\lambda)=0,
 \qquad
 \sum_\lambda(\nu*d)(\lambda)\lambda\lambda^{\mathsf T}
       =\left(\sum_\beta d(\beta)\right)C.
\end{equation}
Indeed, substitute $\lambda=\beta+\gamma$ in the double sum. The terms
in $\beta\beta^{\mathsf T}$, $\beta\gamma^{\mathsf T}$ and
$\gamma\beta^{\mathsf T}$ vanish by the zeroth and first moments of
$\nu$. Thus, with
\[
 \zeta_n=n^{-2}\sum_\lambda d_n(\lambda)
        =n^{-k}\sum_\lambda b(\lambda/n),
\]
the identity in \eqref{pairs:exact-moment} holds exactly for
$(\mu_n,\zeta_n)$, and Riemann sums give $\zeta_n\to\int b>0$.

For $x=\lambda/n$, Taylor expansion gives
\begin{align*}
 n^k\mu_n(\lambda)
 &=n^2\sum_\gamma\nu(\gamma)b(x-\gamma/n)\\
 &=\frac12\sum_{j,l}C_{jl}\partial_j\partial_l b(x)+O(n^{-1})
  =D_Cb(x)+O(n^{-1}).
\end{align*}
The error is uniform in $x$: the third derivatives of $b$ are bounded,
and $\sum_\gamma|\nu(\gamma)|\|\gamma\|^3$ is finite. Moreover, the
supports of both $\mu_n(\lambda)$ and $D_Cb(\lambda/n)$, after division
of the lattice by $n$, lie in a single fixed ball for all $n\geq1$.
There are $O(n^k)$ relevant lattice points and the weight $\|Ux\|^2$
is bounded on that ball. The $1$-Lipschitz property of the negative part
therefore shows that
\[
 \sum_\lambda\mu_n(\lambda)_-\|U\lambda/n\|^2
 -n^{-k}\sum_\lambda(D_Cb(\lambda/n))_-\|U\lambda/n\|^2
 =O(n^{-1}).
\]
The second sum is a Riemann sum for a continuous compactly supported
function. By \eqref{pairs:bump-bound}, it follows that
\[
 \sum_\lambda\mu_n(\lambda)_-\|U\lambda/n\|^2
 \longrightarrow\int(D_Cb)_-\|Ux\|^2
 <\eta\int b.
\]
Fix a sufficiently large $n$ such that $\zeta_n>0$ and
\eqref{pairs:small-negative} holds strictly for $(\mu_n,\zeta_n)$.

It remains to obtain rational coefficients. Let $E\subset\Gamma$ be
a finite set supporting $d_n$. For functions $d$ supported in $E$, the
maps
\[
 d\longmapsto\nu*d,
 \qquad d\longmapsto\zeta=n^{-2}\sum_\lambda d(\lambda)
\]
are linear, and the function
\[
 d\longmapsto
 \sum_{\lambda\in\supp\nu+E}(\nu*d)(\lambda)_-\|U\lambda/n\|^2
\]
is continuous on $\R^E$. The two strict inequalities $\zeta>0$ and
\eqref{pairs:small-negative} consequently persist in an open neighborhood
of $d_n$. Choose a rational vector $d$ in that neighborhood and extend it
by zero outside $E$. Put $\mu=\nu*d$. Both $\mu$ and $\zeta$ are rational.
Convolution still gives zero quotient pushforwards, and
\eqref{pairs:convolution-moments} still gives the exact second moment
\eqref{pairs:exact-moment}. This proves the lemma.
\end{proof}

We have now obtained finite rational weights with exact cancellation
and the required second moment. The inequalities will only be used to
add a common positive semidefinite Gram increment; no further
approximation is needed.

\begin{proof}[Proof of Lemma~\ref{pairs:two-scale}]
Use the notation preceding Lemma~\ref{stencil:finite} and the function
$\nu$ supplied by that lemma:
$w_j=(v_j,u_j)$, $w(\lambda)=\sum_j\lambda_jw_j$,
$U=(u_1\ \cdots\ u_k)$, $\Gamma=\Z^k$, and
\[
 \Lambda_i=\{\lambda\in\Gamma:e_i(w(\lambda))=0\}.
\]
The rational symmetric matrix $C$ of the tensor satisfies
$UCU^{\mathsf T}=I_d$. Apply Lemma~\ref{pairs:rational} to obtain
$n,\mu,\zeta$. Choose a positive integer $M$ such that every
$M\mu(\lambda)$ is an integer. Form two finite lists of affine rows:
the positive list has $M\mu(\lambda)_+$ repetitions of $w(\lambda/n)$,
and the negative list has $M\mu(\lambda)_-$ repetitions.

The zero total mass of $\mu$ shows that the two lists have the same
length. For each label $i$,
\[
 e_i(w(\lambda/n))=e_i(w(\lambda'/n))
 \quad\Longleftrightarrow\quad \lambda-\lambda'\in\Lambda_i.
\]
Thus the zero pushforward on $\Gamma/\Lambda_i$ says exactly that,
after evaluating at $a_i$, every real coordinate value appears equally
often in the two lists. The resulting evaluation vectors are coordinate
permutations of one another.

Divide every affine row in both lists by $\sqrt{\zeta M}$. Their
gradient Gram matrices become
\[
 B_\pm=\frac1\zeta\sum_\lambda\mu(\lambda)_\pm
       (U\lambda/n)(U\lambda/n)^{\mathsf T}.
\]
Equations~\eqref{pairs:exact-moment} and
\eqref{pairs:small-negative} give
\[
 B_+-B_-=I_d,
 \qquad \tr(B_-)<\eta.
\]
Since $B_-$ is positive semidefinite, each of its eigenvalues is at most
its trace. Hence $P=\eta I_d-B_-$ is positive semidefinite. Choose
vectors $r_1,\ldots,r_d\in\R^d$ with
$\sum_jr_jr_j^{\mathsf T}=P$, for example the columns of $P^{1/2}$,
and append the same affine rows $(0,r_j)$ to both lists. The two Gram
matrices are now $(1+\eta)I_d$ and $\eta I_d$. Identical appended rows
preserve equality of all coordinate multisets.

Finally divide every row in both lists by $\sqrt{1+\eta}$, and let
$f_i,g_i$ be their respective evaluation vectors at $a_i$. Their
coordinate multisets remain equal. For any affine row $(v,u)$, the
difference of its evaluations at $a_i$ and $a_j$ is
$(a_i-a_j)\cdot u$. The two gradient Gram identities therefore give
\eqref{pairs:distances} exactly, with $q=\eta/(1+\eta)$.
\end{proof}
